\documentclass[a4paper,11pt]{article}
\usepackage[T2A]{fontenc}
\usepackage[utf8]{inputenc}
\usepackage[russian]{babel}
\usepackage{amsmath,amssymb,mathtools}
\usepackage{geometry}
\usepackage{microtype}
\usepackage{booktabs,tabularx,array}
\usepackage{enumitem}
\usepackage{fancyhdr}
\usepackage{tikz}
\usepackage{xcolor}
\usepackage{hyperref}
\usepackage{parskip}
\geometry{left=24mm,right=24mm,top=25mm,bottom=25mm}
\definecolor{accent}{HTML}{087A79}
\definecolor{softdark}{HTML}{263B3B}
\definecolor{alert}{HTML}{A14343}
\definecolor{muted}{HTML}{657577}
\hypersetup{colorlinks=true,linkcolor=accent,unicode=true,pdfauthor={Лёвин Артём Александрович},pdftitle={Ревью домашней работы — Марк Гукин — 29.10.2026}}
\setlength{\parindent}{0pt}
\setlength{\parskip}{7pt}
\setlength{\headheight}{22pt}
\setlength{\tabcolsep}{5pt}
\renewcommand{\arraystretch}{1.28}
\pagestyle{fancy}
\fancyhf{}
\fancyhead[L]{\footnotesize\color{muted} Тригонометрия: формулы и значения}
\fancyhead[R]{\footnotesize\color{muted} Марк Гукин}
\renewcommand{\headrulewidth}{0.35pt}
\renewcommand{\headrule}{\hbox to\headwidth{\color{accent}\leaders\hrule height \headrulewidth\hfill}}
\fancyfoot[C]{\footnotesize\color{muted}\thepage}
\newcommand{\heading}[1]{\par\vspace{8pt}{\large\color{accent}\bfseries #1}\par\vspace{2pt}}
\newcommand{\smallheading}[1]{\par\vspace{4pt}{\color{softdark}\bfseries #1}\par\vspace{-1pt}}
\begin{document}
\pagenumbering{Roman}
\thispagestyle{empty}
\begin{center}
\vspace*{22mm}
{\small\color{accent} МАТЕМАТИКА}\par\vspace{18mm}
{\LARGE\color{softdark} Ревью домашней работы}\par\vspace{8mm}
{\Large\color{accent} Тригонометрия}\par\vspace{4mm}
{\normalsize Формулы приведения и двойного угла}\par
\vfill
{\normalsize Марк Гукин}\par\vspace{13mm}
{\normalsize 29 октября 2026 г.}\par\vspace{14mm}
{\small Лёвин Артём Александрович}\par\vspace{2mm}
{\small эксклюзивно для Марка Гукина}\par
\vspace{14mm}
\begin{tikzpicture}[x=1cm,y=1cm]
\draw[accent,line width=0.8pt] (0,0) .. controls (1.0,.42) and (1.6,-.20) .. (3,0)
 .. controls (4.1,.27) and (4.7,-.39) .. (6,0);
\end{tikzpicture}
\end{center}
\clearpage
\pagenumbering{arabic}
\heading{Сводная таблица}
\vspace{6pt}
Проверка выполнена по условию и письменному ходу решения. В таблице оставлена только запись, требующая исправления.
\par\vspace{9pt}
\small
\begin{tabularx}{\linewidth}{@{}p{10mm}p{26mm}X p{21mm}p{23mm}@{}}
\toprule
\textbf{№} & \textbf{Тема} & \textbf{Суть ошибки} & \textbf{Ваш ответ} & \shortstack[l]{\textbf{Верный}\\\textbf{ответ}}\\
\midrule
\hyperref[task:8]{\textbf{8}} & Двойной угол & Пропущены скобки после множителя 18; из-за этого промежуточные равенства неверны. & $14$ & $14$\\
\bottomrule
\end{tabularx}
\par\normalsize
\vspace{9pt}
\textcolor{accent}{\textbf{Обратите внимание.}} Совпадение итогового ответа с правильным значением здесь случайно с точки зрения записанной цепочки: последний переход использует скобки, которых нет в предыдущих строках.
\clearpage
\thispagestyle{empty}
\null\vfill
\begin{center}
{\LARGE\color{accent} Разбор ошибок}\\[13pt]
{\normalsize\color{muted} Как правильно раскрывать скобки при подстановке формулы двойного угла}
\end{center}
\vfill\null
\clearpage
\phantomsection\label{task:8}
\heading{Задание 8. Формула двойного угла}
\smallheading{Условие}
Найдите $18\cos 2\alpha$, если $\sin\alpha=-\dfrac13$.
\smallheading{Ваше решение}
После выбора нужной формулы записано:
\[
18\cdot 1-2\sin^2\alpha,
\qquad
18\cdot1-2\cdot\frac19,
\qquad
18\cdot1-\frac29.
\]
Затем произведено вычисление $18\cdot\dfrac79$ и получен ответ $14$.

\smallheading{Ваш ответ}
$14$.

\smallheading{Ошибка}
\textcolor{alert}{\textbf{Ошибка:}} при замене $\cos2\alpha$ выражением $1-2\sin^2\alpha$ потеряны скобки вокруг всей разности.

\smallheading{Где именно возникает ошибка}
До применения формулы исходное выражение $18\cos2\alpha$ записано правильно. Формула двойного угла тоже выбрана верно:
\[
\cos2\alpha=1-2\sin^2\alpha.
\]
Первый некорректный переход возникает, когда вместо
\[
\textcolor{accent}{18\bigl(1-2\sin^2\alpha\bigr)}
\]
появляется
\[
\textcolor{alert}{18\cdot1-2\sin^2\alpha}.
\]
Вторая запись означает, что множитель $18$ относится только к единице.

\smallheading{Почему так нельзя}
При подстановке выражения вместо одного множителя весь подставляемый фрагмент нужно заключать в скобки. По распределительному закону
\[
18(1-2t)=18-36t,
\qquad
18\cdot1-2t=18-2t.
\]
Это разные выражения. В частности, при $t=\dfrac19$:
\[
18\cdot1-\frac29=\frac{160}{9}\ne14.
\]
Поэтому перейти от написанной Вами строки к $18\cdot\dfrac79$ нельзя. Для такого перехода сначала необходимо вычислить разность \emph{в скобках}.

\smallheading{Ключевая идея}
\textcolor{accent}{\textbf{Ключевая идея:}} сначала полностью подставляем формулу двойного угла со скобками, затем находим значение внутри скобок и только после этого умножаем на $18$.

\smallheading{Исправление от последнего правильного шага}
\[
18\cos2\alpha
=18\bigl(1-2\sin^2\alpha\bigr).
\]
\clearpage
\heading{Задание 8. Продолжение разбора}
\smallheading{Полное правильное решение}
По условию $\sin\alpha=-\dfrac13$. Возводим обе части в квадрат:
\[
\sin^2\alpha=\left(-\frac13\right)^2=\frac19.
\]
Применяем формулу двойного угла и сохраняем скобки:
\begin{align*}
18\cos2\alpha
&=18\bigl(1-2\sin^2\alpha\bigr)\\
&=18\left(1-2\cdot\frac19\right)\\
&=18\left(1-\frac29\right)\\
&=18\cdot\frac79\\
&=\frac{18}{9}\cdot7\\
&=2\cdot7=14.
\end{align*}
Здесь $18$ умножает на всю разность $1-\dfrac29$, поскольку эта разность находится в скобках.

\smallheading{Проверка другим допустимым способом}
Можно раскрыть скобки по распределительному закону:
\[
18\left(1-2\cdot\frac19\right)
=18-36\cdot\frac19
=18-4
=14.
\]
Оба способа дают один результат. Ограничений на $\alpha$ при использовании приведённой формулы нет.

\smallheading{Итог}
\textcolor{accent}{\textbf{Итог:}} $14$. Итоговое число у Вас правильное; необходимо исправить запись промежуточных преобразований и сохранять скобки при подстановке формул.

\vspace{8pt}
\smallheading{Задача на закрепление}
Найдите $20\cos2\beta$, если $\sin\beta=\dfrac35$.

\vfill
\begin{center}
{\footnotesize\color{muted} Лёвин Артём Александрович эксклюзивно для Марка Гукина}
\end{center}
\end{document}
